\documentclass[12pt]{article}
\usepackage[utf8]{inputenc}
\usepackage{hyperref}
\usepackage{geometry}
\geometry{margin=1in}

\title{\textbf{Análisis Exploratorio del Rendimiento Estudiantil}}
\author{David Fonseca Lara \\ Juan Esteban Avendaño Castaño}
\date{}

\begin{document}

\maketitle

\section{Introducción}
El presente documento resume el análisis exploratorio realizado en el notebook 
\texttt{EDA and cleaning.ipynb}, el cual utiliza el dataset \textit{Students Performance in Exams} 
de Kaggle. Este dataset contiene información sobre el rendimiento académico de estudiantes en 
tres áreas principales: matemáticas, lectura y escritura.  

El objetivo principal es comprender las desigualdades en el desempeño escolar, 
identificando patrones de distribución, variabilidad y la presencia de casos extremos.

\section{Problema abordado}
Se busca responder preguntas clave como:
\begin{itemize}
    \item ¿Existen diferencias sistemáticas entre el desempeño en matemáticas, lectura y escritura?
    \item ¿Qué tan grande es la variabilidad en los puntajes entre estudiantes?
    \item ¿Se observan valores atípicos que podrían indicar problemas estructurales o anomalías en la recolección de datos?
\end{itemize}

\section{Problemas similares en la literatura}
El tema del rendimiento escolar ha sido ampliamente estudiado, con énfasis en factores socioeconómicos, 
culturales y de género.  

El \textit{Coleman Report} (1966) mostró que el contexto socioeconómico es determinante en los logros académicos, 
más allá de los recursos escolares disponibles.  

De manera complementaria, los informes de la OCDE, en particular el 
\textit{Programme for International Student Assessment (PISA)} desde el año 2000, han documentado brechas 
persistentes en el rendimiento escolar vinculadas al género y al nivel socioeconómico.  

Un metaanálisis de Sirin (2005) confirma que el nivel socioeconómico de los estudiantes y sus familias 
es uno de los predictores más sólidos de éxito académico.

\section{Análisis general del notebook}
El análisis exploratorio realizado ofrece varias observaciones clave:

\begin{enumerate}
    \item \textbf{Distribuciones de puntajes:} Los resultados en matemáticas, lectura y escritura son relativamente balanceados, con medias cercanas a las medianas. Esto sugiere ausencia de sesgos muy fuertes.
    \item \textbf{Variabilidad significativa:} La desviación estándar en torno a 15 puntos refleja una dispersión considerable. Algunos estudiantes alcanzan puntajes sobresalientes (100), mientras otros presentan un rendimiento muy bajo (cercano a 0).
    \item \textbf{Asimetría en lectura y escritura:} Las distribuciones muestran un sesgo hacia valores bajos, evidenciando un grupo de estudiantes con dificultades en estas áreas.
    \item \textbf{Detección de outliers:} Casos como un puntaje de 0 en matemáticas sugieren posibles anomalías: ausencia en el examen, errores de captura de datos o situaciones de vulnerabilidad extrema en el aprendizaje.
\end{enumerate}

En conjunto, los hallazgos muestran un núcleo de estudiantes con rendimiento medio-alto, 
pero también desigualdades profundas que hacen evidente la necesidad de políticas educativas 
enfocadas en atender tanto a los estudiantes en riesgo como en la identificación temprana de 
patrones de desigualdad académica.

\section{Referencia del código}
El código utilizado para el análisis exploratorio se encuentra disponible en el siguiente repositorio:  

\url{https://github.com/Adrephos/machinelearning/blob/master/project/entrega\%201/EDA\_and\_cleaning.ipynb}

\section*{Referencias}
\begin{itemize}
    \item Coleman, J. S., et al. (1966). \textit{Equality of Educational Opportunity}. Washington, DC: U.S. Government Printing Office.
    \item OECD. (2019). \textit{PISA 2018 Results (Volume I): What Students Know and Can Do}. OECD Publishing. \url{https://doi.org/10.1787/5f07c754-en}
    \item Sirin, S. R. (2005). Socioeconomic Status and Academic Achievement: A Meta-Analytic Review. \textit{Review of Educational Research}, 75(3), 417–453. \url{https://doi.org/10.3102/00346543075003417}
\end{itemize}

\end{document}
